\section{Семинар \thesection}
\subsection{Тождественные частицы.}
Задачи по изучению поведения одной частицы в квантовой механике представляют интерес в первую очередь как задачи, которые можно решить аналитически. Однако, в практических применениях гораздо чаще встречаются многочастичные задачи. В этой главе мы начнём рассматривать именно такие задачи.

Обсудим особенности работы с многочастичными задачами в квантовой механике по сравнению с классической. В классической механике для описания нескольких систем мы рассматриваем декартово произведение их фазовых пространств. Другими словами, размерность комплексной системы, состоящей из подсистем, является суммой размерностей этих подсистем. В квантовой механике мы работаем с векторными пространствами, суперпозицией и запутанностью. Эти особенности не позволяют нам использовать декартово произведение для описания сложных квантовых систем. Вместо этого мы будем использовать тензорное произведение $\otimes$. Так, для описания состояния системы, состоящей из двух частиц $\ket{\psi_1}$ и $\ket{\psi_2}$ используется вектор $\ket{\Psi} = \ket{\psi_1} \otimes \ket{\psi_2}$.

Часто, для упрощения записи, знак тензорного произведения будет опускаться. Например, вместо $\ket{\psi_1} \otimes \ket{\psi_2}$ могут писать $\ket{\psi_1}\ket{\psi_2}$. С такой же частотой встречается и нотация вида $\ket{\psi_1\psi_2}$ когда понятно, что подразумевается именно система, состоящая из двух состояний. Аналогичная ситуация при рассмотрении операторов. Так, оператор сложной системы, состоящий из двух операторов, действующих на её подсистемах, запишется как: $\hat{A} = \hat{a}_1 +\hat{a}_2$. Правильной нотацией в данном случае было бы $\hat{A} = \hat{a}_1 \otimes \hat{1}_2 + \hat{1}_1 \otimes \hat{a}_2$. Однако тривиальную часть (единичный оператор) чаще всего опускают.

В дальнейшем я буду пользоваться указанной выше нотацией и не буду ставить шляпки над операторами, если из контекста будет понятно, что речь идёт именно о операторах.

\subsubsection{Сумма моментов.}

Из предыдущего семестра мы узнали, что общий момент $\mathbf{J}$ для частицы состоит из суммы двух компонент -- момента импульса $\mathbf{L}$ и собственного момента $\mathbf{S}$, то есть $\mathbf{J} = \mathbf{L} + \mathbf{S}$. Пришло время раскрыть более подробно свойства оператора собственного момента (или спинового оператора) и какие особенности из этих свойств вытекают. 

Начнём с того, как правильно складывать моменты, если в нашей системе несколько частиц. Необходимость в этом появляется из-за полного набора операторов для описания такой многочастичной системы. Напомню, что для описания сферический-симметричной задачи мы использовали собственные состояния оператора $\mathbf{L}^2$ и $L_z$, а именно состояния $\ket{n,l,m}$. Если учитывать спиновую часть, то должны добавиться собственные состояния $\ket{s, m_s}$ для операторов $\mathbf{S}^2$ и $S_z$. Получается, что полный набор составляют операторы $\mathbf{L}^2$, $\mathbf{S}^2$, $L_z$ и $S_z$. 

В случае задачи многих частиц, эти операторы будут представлять из себя сумму операторов для каждой частицы. Исходя из этого, нам нужно научиться складывать моменты разных частиц. Под складыванием моментов обычно подразумевается переход от собственных состояний операторов отдельных частиц к собственным состояниям их суммы. 

Рассмотрим сначала конкретный пример: спиновую часть частиц со спином $1/2$. Для вычисления суммы спиновых операторов нужно перейти от $S_{1z}$ и $S_{2z}$ к $\mathbf{S}^2 = (\mathbf{S}_1 + \mathbf{S}_2)^2$ и $S_z = S_{1z} + S_{2z}$. Собственные состояния вторых уже были выписаны выше, собственные состояния первых обозначим как $\ket{+}$ и $\ket{-}$, где плюс соответствует проекции $m_s = 1/2$, а минус -- $m_s = -1/2$.

Найдём коэффициенты перехода от одного базиса к другому. Для этого воспользуемся понижающим оператором $S_- = S_{1-} + S_{2-}$. Напомню, как действуют эти операторы:
\begin{gather*}
S_{\pm}\ket{s, m_s} = \sqrt{(s\mp m_s)(s\pm m_s + 1)}\ket{s, m_s - 1};\\
S_{j-}\ket{+} = \sqrt{(1/2+1/2)(1/2 - 1/2 + 1)}\ket{-} = \ket{-},\; S_{j-}\ket{-} = 0;\\
S_{j+}\ket{-} = \sqrt{(1/2 + 1/2)(1/2 - 1/2 + 1)}\ket{+} = \ket{+},\; S_{j+}\ket{+} = 0.
\end{gather*}
Применим понижающий оператор к состоянию $\ket{s=1, m_s=1}$. Понятно, что это состояние отвечает двум спинам с положительной проекцией $m_{1,2} = 1/2$, то есть $\ket{++}$.
\[
S_- \ket{1, 1} = (S_{1-} + S_{2-})\ket{++}\implies \sqrt{2}\ket{1, 0} = \ket{-+} + \ket{+-}
\]
Если применить этот оператор ещё раз, получим
\[
S_- \ket{1, 0} = \frac{1}{\sqrt{2}}(S_{1-} + S_{2-})(\ket{-+}+\ket{+-})\implies \ket{1, -1} = \ket{--}
\]
Для случая $\ket{s=0, m=0}$ нужно построить ортогональный вектор ко всем трём векторам найденным выше. Легко проверить, что вектор $\sqrt{2}\ket{s=0, m=0} = (\ket{+-} - \ket{-+})$ идеально для этого подходит.

Заметим, что мы построили матрицу перехода от базиса $\{m_1, m_2\}$ к базису $\{s, m\}$. Действительно, если выбрать стандартный базис $\ket{+} = \begin{pmatrix}1 & 0\end{pmatrix}^T$, тогда матрица перехода размером $4\times 4$ (так как вектор $\ket{++}$ получен тензорным произведением двух векторов размерности два) будет выглядеть следующим образом
\[
\ket{s, m} = \begin{pmatrix}
1 & 0 & 0 & 0\\
0 & 1/\sqrt{2} & 1/\sqrt{2} & 0\\
0 & 1/\sqrt{2} & -1/\sqrt{2} & 0\\
0 & 0 & 0 & 1\\
\end{pmatrix} \ket{m_1, m_2}
\]
Получается, чтобы определить, как складывать моменты в многочастичной системе, необходимо найти матрицу перехода от базиса моментов двух части к их общему базису. Мы сделали это для конкретного случая, когда спин частиц равен 1/2. Однако, существует общий подход для определения матрицы перехода. Для того, чтобы подступиться к нему, отмечу, что элементы полученной матрицы являются частным примером \textit{коэффициентов Клебша-Гордона}. Как раз умение находить такие коэффициенты и даст нам общую "формулу" для сложения моментов разных частиц.

Попробуем вывести общий подход к нахождению коэффициентов Клебша-Гордона. Для этого рассмотрим случай, когда у нас есть два произвольных угловых момента $\mathbf{J}_1$ и $\mathbf{J}_2$. Из прошлого семестра известно, что полный набор операторов для такой системы будет состоять из $\mathbf{J}^2_1$, $\mathbf{J}^2_2$, $J_{1z}^2$ и $J_{2z}^2$. Этому набору операторов будет отвечать следующий базис:
\begin{gather*}
\mathbf{J}^2_1\ket{j_1j_2;m_1m_2} = j_1(j_1+1)\hbar^2\ket{j_1j_2; m_1m_2}\\
\mathbf{J}^2_2\ket{j_1j_2;m_1m_2} = j_2(j_2+1)\hbar^2\ket{j_1j_2; m_1m_2}\\
J_{z1}\ket{j_1j_2;m_1m_2} = m_1\hbar\ket{j_1j_2; m_1m_2}\\
J_{z2}\ket{j_1j_2;m_1m_2} = m_1\hbar\ket{j_1j_2; m_1m_2}
\end{gather*}
От них, как и в рассмотренном выше примере, необходимо перейти к операторам $\mathbf{J}^2$, $\mathbf{J}^2_1$, $\mathbf{J}^2_2$ и $J_z$, где ненумерованные операторы отвечают сумме моментов. Базис новых операторов:
\begin{gather*}
\mathbf{J}^2_1\ket{j_1j_2;jm} = j_1(j_1+1)\hbar^2\ket{j_1j_2; jm}\\
\mathbf{J}^2_2\ket{j_1j_2;jm} = j_2(j_2+1)\hbar^2\ket{j_1j_2; jm}\\
\mathbf{J}^2\ket{j_1j_2;jm} = j(j+1)\hbar^2\ket{j_1j_2; jm}\\
J_{z}\ket{j_1j_2;jm} = m\hbar\ket{j_1j_2; jm}
\end{gather*}

Для перехода от одного базиса к другому воспользуемся стандартным свойством базиса:
\[
\sum_{m_1}\sum_{m_2}\ket{j_1j_2;m_1m_2}\bra{j_1j_2;m_1m_2}=1
\]
Тогда, новый базис можно записать в следующем виде:
\[
\ket{j_1j_2;jm}=\sum_{m_2}\ket{j_1j_2;m_1m_2}\braket{j_1j_2;m_1m_2}{j_1j_2;jm}
\]
Отсюда видно, что $\braket{j_1j_2;m_1m_2}{j_1j_2;jm}$ и будут теми самыми коэффициентами Клебша-Гордона. 
\subsubsection{Фермионы и Бозоны}
